\chapter{Conclusion}\label{ch:concl}

\section{Summary of Contributions}

This work presented a complete machine learning system for predicting the
performance of photoelectrochemical water-splitting devices.
The main contributions are:

\begin{enumerate}
  \item \textbf{Curated and augmented PEC dataset.}
        A dataset of 676 entries was assembled from 529 published experimental
        results and extended through physics-based estimation of STH and
        H\textsubscript{2} values from measured photocurrents (Faraday's law /
        STH formula) and through 147 literature-anchored synthetic rows.
        The augmentation increased STH coverage from 9\% to 82\% and
        H\textsubscript{2} coverage from 11\% to 83\%.

  \item \textbf{Physics-informed feature engineering.}
        Twenty numerical features were engineered from eight raw inputs,
        including interaction terms motivated by the Nernst equation
        (bandgap $\times$ pH), Butler--Volmer kinetics (bias $\times$ ionic strength),
        and Arrhenius-like thermal activation (pH $\times$ temperature).

  \item \textbf{Ensemble ML pipeline with uncertainty quantification.}
        XGBoost and Random Forest regressors are trained per target with
        5-fold CV.
        An ensemble inference strategy weights predictions by CV R\textsuperscript{2}
        and computes a confidence score from model agreement (Equation~\ref{eq:confidence}).

  \item \textbf{Data-driven performance classification.}
        A composite score (Equation~\ref{eq:score}) calibrated to training
        percentiles and modulated by model confidence replaces the previous
        hardcoded threshold, enabling the HIGH class to be predicted for
        genuinely exceptional inputs.

  \item \textbf{Deployed full-stack application.}
        A Flask REST API and single-page web interface expose the models
        with real-time predictions, confidence scores, and interactive
        dataset exploration.

  \item \textbf{Measurable model improvement through augmentation.}
        CV R\textsuperscript{2} improved from 0.065 to 0.454 for Photocurrent,
        from 0.278 to 0.533 for STH, and from $-$0.090 to 0.366 for
        H\textsubscript{2} after augmentation and feature expansion.
\end{enumerate}

\section{Limitations and Future Work}

Several directions remain open for future investigation:

\begin{itemize}
  \item \textbf{Real data collection.}
        The most impactful improvement would be expanding the number of
        genuine experimental STH and H\textsubscript{2} measurements.
        Scraping open-access papers with NLP-based table extraction tools
        (e.g., GROBID, PaperScraper) could double the real dataset within
        a semester.

  \item \textbf{DFT-derived features.}
        Incorporating computed properties (formation energy, effective masses,
        dielectric constant) from databases such as the Materials Project
        \cite{jain2013commentary} or AFLOW \cite{curtarolo2012aflow} would
        add material-specific information not captured by composition alone.

  \item \textbf{Advanced models.}
        For datasets exceeding 2000 entries, Bayesian neural networks or
        deep ensembles with Monte Carlo dropout would provide calibrated
        uncertainty intervals.
        Graph neural networks acting on crystal structures are a longer-term
        option when structural data becomes available.

  \item \textbf{Multi-target learning.}
        A joint model predicting $J_{ph}$, STH, and HER simultaneously
        could exploit the correlations between targets and reduce
        data requirements.

  \item \textbf{Active learning loop.}
        The confidence score could drive an active learning workflow,
        identifying which experimental configurations would most reduce
        model uncertainty and recommending them to researchers.
\end{itemize}

Despite its current limitations, the system provides a practical and
immediately usable tool for accelerating the pre-screening of PEC
photoelectrode candidates, reducing the number of synthesise-and-test
iterations required before finding high-efficiency configurations.
